% summary.tex
% mainfile: perfbook.tex
% SPDX-License-Identifier: CC-BY-SA-3.0

\chapter{回顾与展望}
\label{chp:Looking Forward and Back}
%
\Epigraph{历史就是本可以避免的事情的总和。}
	  {Konrad Adenauer}

你已经读完了这本书，干得好！
我希望你的旅程是愉快的，同时也是充满挑战和有价值的。

对于编辑和贡献者来说，这是通向第二版的旅程的终点，
但对于愿意加入的人来说，这也是通向第三版的旅程的起点。
无论如何，回顾这段旅程是件好事。

\Cref{chp:How To Use This Book}介绍了本书的内容，
以及为那些对底层并行编程以外的内容感兴趣的读者提供的一些替代选择。

\Cref{chp:Introduction}涵盖了并行编程的挑战以及应对这些挑战的
高层次方法。
它还触及了在避免这些挑战的同时仍然获得并行性大部分好处的途径。

\Cref{chp:Hardware and its Habits}对多核硬件进行了高层次概述，
尤其是那些对并发软件构成挑战的方面。
这一\lcnamecref{chp:Hardware and its Habits}将这些挑战的责任
归于其应得之处——很大程度上归于物理定律，
而较少归于固执的硬件架构师和设计师。
然而，硬件架构师和工程师可能能做一些事情，
这一\lcnamecref{chp:Hardware and its Habits}讨论了其中一些。
与此同时，软件架构师和工程师必须尽其本分来应对这些挑战，
正如本书其余部分所讨论的那样。

\Cref{chp:Tools of the Trade}
快速概述了底层并发编程的常用工具。
\Cref{chp:Counting}随后展示了这些工具的使用——
更重要的是，展示了并行编程设计技术的使用——
应用于简单但出人意料地具有挑战性的并发计数任务。
事实上如此具有挑战性，以至于有多种并发计数算法被广泛使用，
每种都专门针对不同的使用场景。

\Cref{chp:Partitioning and Synchronization Design}更深入地探讨了
最重要的并行编程设计技术，即在尽可能高的层次上对问题进行分区。
这一\lcnamecref{chp:Partitioning and Synchronization Design}还
概述了这一设计空间中的若干要点。

\Cref{chp:Locking}详细阐述了并行编程的主力（也是恶棍）——locking。
这一\lcnamecref{chp:Locking}涵盖了多种类型的 locking，
并针对 locking 的许多众所周知且被大力宣传的缺点
提出了一些工程解决方案。

\Cref{chp:Data Ownership}讨论了 data ownership 的使用，
其中同步由给定数据项与特定 thread 的关联来提供。
在适用的情况下，这种方法将出色的性能和可扩展性与
深刻的简洁性结合在一起。

\Cref{chp:Deferred Processing}展示了一点点拖延如何能够
大大提高性能和可扩展性，同时在出人意料的大量情况下
还能简化代码。
这一\lcnamecref{chp:Deferred Processing}中介绍的多种机制
利用了 CPU cache 复制只读数据的能力，
从而绕开了残酷地限制光速和原子大小的物理定律。
\Cref{chp:Data Structures}考察了并发数据结构，
重点是 hash table，它在并行程序中有着悠久而光荣的历史。

\Cref{chp:Validation}深入研究了代码审查和测试方法，
\cref{chp:Formal Verification}概述了形式化验证。
无论你站在形式化验证/测试之争的哪一边，
如果代码没有经过彻底验证，它就不能正常工作。
对于并发代码来说，这一点至少加倍适用。

\Cref{chp:Putting It All Together}展示了若干场景，
其中将并发机制相互组合或与其他设计技巧组合，
可以极大地改善并行程序员的工作。
\Cref{sec:advsync:Advanced Synchronization}考察了高级同步方法，
包括无锁编程、\IXacrl{nbs}和并行实时计算。
\Cref{chp:Advanced Synchronization: Memory Ordering}深入探讨了
至关重要的 memory ordering 主题，
展示了帮助你不仅解决 memory ordering 问题，
而且完全避免这些问题的技术和工具。
\Cref{chp:Ease of Use}简要概述了易用性这一
出人意料地重要的主题。

最后但绝非最不重要的是，\cref{chp:Conflicting Visions of the Future}
详细阐述了若干关于未来的矛盾愿景，
包括 CPU 技术趋势、transactional memory、
hardware transactional memory、在回归测试中使用形式化验证，
以及长久以来关于并行编程的未来属于函数式编程语言的预测。

但现在我们已经回顾了第二版的内容，
这本书是如何开始的呢？

Paul 的并行编程之旅真正开始于 1990 年，
那年他加入了 Sequent Computer Systems, Inc.。
Sequent 采用类似学徒制的培训方式，
新聘工程师被安排在经验丰富的工程师的隔间旁边，
后者指导他们、审查他们的代码，
并在各种主题上给予大量建议。
一些新聘工程师极大地受益于当时没有片上 cache 的事实，
这意味着逻辑分析仪可以轻松显示给定 CPU 的指令流和内存访问，
附带精确的时序信息。
当然，这种透明性的代价是 CPU 核心时钟频率
比 21 世纪的慢 100 倍。
凭借学徒制和硬件性能的透明性，
这些新聘工程师在两三个月内就成为了高效的并行程序员，
其中一些人在几年内就做出了开创性的工作。

Sequent 深知其快速培训新工程师掌握并行性奥秘的能力是非凡的，
因此出版了一本精炼的小册子，
结晶了公司的并行编程智慧~\cite{SQNTParallel}，
与几年前写的两篇开创性论文~\cite{Beck85,Inman85}相得益彰。
已经精通这些奥秘的人对这本书和这些论文表示赞赏，
但新手通常无法从中获益太多，
总是犯出极具创意但极具破坏性的错误，
而这些错误在书中和论文中都没有被明确禁止。\footnote{
	``但你到底为什么要\emph{那样}做？？？''
	``呃，为什么不呢？''}
这种情况当然使 Paul 开始考虑写一本改进的书，
但在这段时间里，他的努力仅限于内部培训材料和已发表的论文。

到 1999 年 Sequent 被 IBM 收购时，
世界上许多最大的数据库实例都运行在 Sequent 硬件上。
但时代在变，到 2001 年，Sequent 的许多并行程序员
已经将注意力转向了 Linux 内核。
在一些最初的犹豫之后，Linux 内核社区热情且有效地
拥抱了并发~\cite{SilasBoydWickizer2010LinuxScales48,McKenney:2012:BEP:2414729.2414734}，
整个社区贡献了许多出色的创新和改进。
Paul 不时想到写一本书，但生活节奏飞快，
所以他在这个项目上没有取得进展。

2006 年，Paul 受邀参加一个关于 Linux 可扩展性的会议，
并获得了在一个由著名并行编程专家组成的小组讨论中
提出最后一个问题的殊荣。
Paul 在提问时指出，在 1991 年到 2006 年的 15 年间，
并行系统的价格从一栋房子的价格降到了一辆中档自行车的价格，
而且很明显，在未来 15 年（直到 2021 年）
还有很大的进一步大幅降价空间。
他还指出，价格下降应该会带来更多的熟悉度，
以及在解决并行编程问题方面更快的进展。
这引出了他的问题：
``到 2021 年，为什么并行编程不会变得稀松平常？''

第一位小组成员似乎相当蔑视任何会问如此荒谬问题的人，
并迅速用一句简短的回答作出回应。
对此 Paul 给出了同样简短的回应。
他们来来回回了一段时间，例如，小组成员简短回答的
``Deadlock''引出了 Paul 简短的回应``Lock dependency checker''。

这位小组成员最终用完了简短的回答，即兴发挥了最后一句：
``像你这样的人应该被锤子敲头！''

Paul 的回应当然是``你得排队等候！''

Paul 将注意力转向了下一位小组成员，后者似乎在
同意第一位小组成员和不想应对 Paul 的一连串回应之间左右为难。
因此他做了一个简短的不置可否的发言。
其余的小组成员也大致如此。

直到轮到最后一位小组成员——一个你可能听说过的人，
名叫 Linus Torvalds。
Linus 指出，三年前（即 2003 年），
任何与并发相关的补丁的初始版本通常质量很差，
存在设计缺陷和许多 bug。
即使清理到足以被接受，bug 仍然存在。
Linus 将此与当时 2006 年的情况进行了对比，他说
并发相关补丁的第一个版本就设计良好、几乎没有甚至没有 bug
的情况并不罕见。
然后他建议，\emph{如果}工具继续改进，
那么\emph{也许}并行编程到 2021 年会变得稀松平常。\footnote{
	工具确实在持续改进，包括模糊测试器、
	lock 依赖检查器、静态分析器、形式化验证、
	内存模型以及 coccinelle 等代码修改工具。
	因此，那些想要断言 2021 年的并行编程并非稀松平常的人，
	应该参考\cref{chp:Introduction}的题词。}

会议随后结束了。
Paul 并不惊讶许多观众对他敬而远之，
尤其是那些与第一位小组成员看法相同的人。
Paul 也不惊讶有一些观众感谢他提出了这个问题。
然而，当一个人走过来泪流满面地说``谢谢你''，
哽咽得几乎说不出话时，他相当惊讶。

你看，这个人曾在 Sequent 工作了几年，
因此非常精通并行编程。
此外，他目前被分配到一个负责编写并行代码的团队中。
但事情进展得并不顺利。
不是他们在理解他对并行编程的解释方面有困难。

而是他们根本\emph{拒绝听他说话}。

简而言之，他的团队对这个人的方式，
就像第一位小组成员试图对待 Paul 的方式一样。
于是在那一刻，Paul 从``我将来应该写一本书''
变成了``我将竭尽全力写这本书''。
Paul 不好意思地承认，他不记得那个人的名字，
甚至可能从未知道过。

尽管如此，这本书是为那个人而写的。

\begin{figure}
\centering
\resizebox{3in}{!}{\includegraphics{cartoons/UseTheRightToolsBubble}}
\caption{最重要的一课}
\ContributedBy{fig:summary:The Most Important Lesson}{Melissa Broussard}
\end{figure}

这本书也是为所有希望将底层并发编程加入其技能组合的人而写的。
如果你只记住了这本书中的一件事，
那就记住\cref{fig:summary:The Most Important Lesson}的教训吧。

这本书也是向那位未具名小组成员的未具名雇主致敬。
几年后，这位雇主选择任命了一位拥有更有用的经验
和更少简短回答的人。
此人也在一个小组讨论中，在那次讨论中，
他直视着我说并行编程也许只比顺序编程难 5\pct。

对于我们其余的人来说，当有人试图向我们展示一个
紧迫问题的解决方案时，
也许我们至少应该给他们一个倾听的礼遇！
